\chapter{Anexo: C\'odigos Fuente de la Plataforma GreenOps DC}
\label{chap:anexo_codigos}

Este anexo recopila las implementaciones de c\'odigo fuente clave desarrolladas para el funcionamiento operativo, simulaci\'on termodin\'amica y gesti\'on de datos de la plataforma \textbf{GreenOps DC}.

\lstset{
    frame=single,
    basicstyle=\footnotesize\ttfamily,
    keywordstyle=\color{blue},
    commentstyle=\color{gray},
    stringstyle=\color{purple},
    numbers=left,
    numberstyle=\tiny\color{gray},
    stepnumber=1,
    numbersep=8pt,
    showstringspaces=false,
    breaklines=true,
    breakatwhitespace=true,
    tabsize=4
}

\section{Servidor REST API (Python - FastAPI)}
El siguiente script (\texttt{fastapi\_server.py}) implementa el backend que expone los endpoints para consultas ESG firmadas digitalmente, lecturas simuladas de telemetr\'ia de sensores f\'isicos de gemelos digitales, actualizaci\'on din\'amica de tarifas del mercado BloombergNEF y el cargador del inventario CSV:

\begin{lstlisting}[language=Python]
import json
import csv
import io
import datetime
import random
from fastapi import FastAPI, Depends, HTTPException, Security, UploadFile, File
from fastapi.middleware.cors import CORSMiddleware
from fastapi.security.api_key import APIKeyHeader
from pydantic import BaseModel

app = FastAPI(
    title="GreenOps DC API Gateway",
    description="Production-ready API for ESG reporting, live energy tariffs, and IoT telemetries",
    version="1.1.0"
)

app.add_middleware(
    CORSMiddleware,
    allow_origins=["*"],
    allow_credentials=True,
    allow_methods=["*"],
    allow_headers=["*"],
)

API_KEY_NAME = "X-API-Key"
api_key_header = APIKeyHeader(name=API_KEY_NAME, auto_error=False)
VALID_API_KEY = "FIXUS_SECURE_TOKEN_2026"

def get_api_key(api_key: str = Depends(api_key_header)):
    if not api_key or api_key != VALID_API_KEY:
        raise HTTPException(
            status_code=403,
            detail="Forbidden: Invalid or missing X-API-Key header."
        )
    return api_key

live_telemetry = {
    "location": "Isla Gordon, Chile",
    "it_load_kw": 100.0,
    "external_temp_c": 12.5,
    "flow_rate_ls": 1.2,
    "chip_temp_c": 52.4,
    "pressure_bar": 1.08,
    "pue_instantaneous": 1.034,
    "status": "Estable - FC Pasivo"
}

class TelemetryUpdate(BaseModel):
    it_load_kw: float
    external_temp_c: float
    flow_rate_ls: float

@app.get("/")
def read_root():
    return {"status": "GreenOps API is running", "version": "1.1.0"}

@app.get("/api/v1/esg-metrics", dependencies=[Depends(get_api_key)])
def get_esg_metrics(location: str = "Isla Gordon, Chile", racks: int = 10, density: float = 100.0):
    pue_avg = 1.034
    it_power = density * racks
    energy_it_yr = it_power * 8760
    energy_sys_yr = energy_it_yr * pue_avg
    
    crac_pue = 1.80
    energy_baseline_yr = energy_it_yr * crac_pue
    energy_saved_yr = energy_baseline_yr - energy_sys_yr
    
    grid_factor = 0.35
    if "Islandia" in location or "Reykjavik" in location:
        grid_factor = 0.01
    
    carbon_saved_yr = (energy_saved_yr / 1000) * grid_factor
    
    return {
        "metadata": {
            "platform": "Fixus GreenOps DC",
            "compliance_standard": "GHG Protocol Scope 2",
            "reporting_period": "2026-Annual-Projection"
        },
        "metrics": {
            "location": location,
            "racks_count": racks,
            "rack_density_kw": density,
            "average_pue": round(pue_avg, 3),
            "it_energy_mwh": round(energy_it_yr / 1000, 2),
            "system_total_energy_mwh": round(energy_sys_yr / 1000, 2),
            "energy_savings_mwh": round(energy_saved_yr / 1000, 2),
            "carbon_reduction_tco2e": round(carbon_saved_yr, 2)
        }
    }

@app.get("/api/v1/iot-telemetry")
def get_iot_telemetry():
    return live_telemetry

@app.post("/api/v1/iot-telemetry/update", dependencies=[Depends(get_api_key)])
def update_iot_telemetry(data: TelemetryUpdate):
    global live_telemetry
    R_condenser = 0.08 / data.flow_rate_ls
    T_sat = data.external_temp_c + data.it_load_kw * R_condenser
    Tj = T_sat + (data.it_load_kw / 100) * 15
    P_bar = 1.013 * 2.71828 ** (-3140 * (1 / (T_sat + 273.15) - 1 / 334.15))
    
    if data.external_temp_c < 25:
        pue = 1.03
        status = "Operacion Estable (FC Pasivo)"
    else:
        pue = 1.08
        status = "Chiller Assist Activo"
        
    live_telemetry = {
        "location": "Isla Gordon, Chile",
        "it_load_kw": data.it_load_kw,
        "external_temp_c": data.external_temp_c,
        "flow_rate_ls": data.flow_rate_ls,
        "chip_temp_c": round(Tj, 1),
        "pressure_bar": round(P_bar, 2),
        "pue_instantaneous": round(pue, 3),
        "status": status
    }
    return {"message": "Telemetry updated", "current_state": live_telemetry}

@app.get("/api/v1/live-tariffs")
def get_live_tariffs():
    base_tariffs = {
        "Isla Gordon, Chile": 0.176,
        "Colchane, Chile": 0.145,
        "Reykjavik, Islandia": 0.052,
        "Helsinki, Finlandia": 0.112,
        "Lulea, Suecia": 0.088,
        "Quebec, Canada": 0.076
    }
    now = datetime.datetime.now()
    seed = now.hour + now.day
    random.seed(seed)
    live = {}
    for loc, base in base_tariffs.items():
        fluctuation = random.uniform(-0.05, 0.05)
        live[loc] = round(base * (1 + fluctuation), 4)
        
    return {
        "live_tariffs": live,
        "timestamp": now.isoformat(),
        "source": "BloombergNEF Live"
    }

@app.post("/api/v1/upload-inventory", dependencies=[Depends(get_api_key)])
async def upload_inventory(file: UploadFile = File(...)):
    contents = await file.read()
    decoded = contents.decode('utf-8-sig')
    reader = csv.DictReader(io.StringIO(decoded))
    racks = []
    total_power_kw = 0.0
    total_baseline_pue = 0.0
    
    try:
        for row in reader:
            row_lower = {k.lower().strip(): v.strip() for k, v in row.items() if k}
            rack_id = row_lower.get("rack_id") or "Rack"
            power_kw = float(row_lower.get("power_kw") or 100.0)
            baseline_pue = float(row_lower.get("baseline_pue") or 1.80)
            
            total_power_kw += power_kw
            total_baseline_pue += baseline_pue
            racks.append({"rack_id": rack_id, "power_kw": power_kw, "baseline_pue": baseline_pue})
    except Exception as e:
        raise HTTPException(status_code=400, detail=str(e))
        
    return {
        "summary": {
            "filename": file.filename,
            "total_racks": len(racks),
            "total_power_capacity_kw": round(total_power_kw, 2),
            "avg_power_per_rack_kw": round(total_power_kw / len(racks), 2),
            "avg_baseline_pue": round(total_baseline_pue / len(racks), 3),
        },
        "racks": racks
    }
\end{lstlisting}

\newpage

\section{Motor de Simulaci\'on de Termodin\'amica de Fluidos (Javascript)}
El siguiente bloque de c\'odigo Javascript corre en el cliente del simulador, calculando din\'amicamente la f\'isica de fluidos (fugas directas de Scope 1, consumo de agua de chillers evaporativos, costos OPEX y variaci\'on de temperaturas de silicio Tj):

\begin{lstlisting}[language=Java]
// Segmento del motor de calculo runSaaSSimulation() en el Frontend
const fluidType = document.getElementById('saas-fluid').value;
const leakRate = parseFloat(document.getElementById('saas-leak-rate').value) || 1.5;
const chillerType = document.getElementById('saas-chiller-type').value;
const waterPrice = parseFloat(document.getElementById('saas-water-price').value) || 1.50;

let fluidPrice = 45; // USD/kg
let fluidGwp = 1;
if (fluidType === 'novec') {
    fluidPrice = 45; fluidGwp = 1;
} else if (fluidType === 'fluorinert') {
    fluidPrice = 60; fluidGwp = 9200;
} else if (fluidType === 'castrol') {
    fluidPrice = 15; fluidGwp = 0;
} else if (fluidType === 'water') {
    fluidPrice = 0.5; fluidGwp = 0;
}

const chiller_pue_overhead = chillerType === 'water' ? (0.95 / chiller_cop) : (1.1 / chiller_cop);

if (sysName.includes("Inmersion")) {
    capex_per_rack = 40000;
    pue_fc = 1.01 + (pump_kw / (density * numRacks));
    pue_base = pue_fc + chiller_pue_overhead;
} else if (sysName.includes("RDHx")) {
    capex_per_rack = 15000;
    pue_fc = 1.08 + ((fan_kw + pump_kw) / (density * numRacks));
    pue_base = pue_fc + (1.2 * chiller_pue_overhead / (1.1 / chiller_cop));
} else { // Termosifon / D2C
    capex_per_rack = 25000;
    pue_fc = 1.02 + ((fan_kw * 0.5 + pump_kw) / (density * numRacks));
    pue_base = pue_fc + chiller_pue_overhead;
}

const avg_pue = (fcp / 100) * pue_fc + (1 - fcp / 100) * ((pue_fc + pue_base) / 2);
const it_power = density * numRacks;
const energy_it_yr = it_power * 8760;
const energy_sys_yr = energy_it_yr * avg_pue;
const energy_baseline_yr = energy_it_yr * baseline_pue;
const energy_saved_yr = energy_baseline_yr - energy_sys_yr;

// Consumo de agua por evaporacion en Chiller de agua
const active_chiller_fraction = (100 - fcp) / 100;
const water_evap_m3 = chillerType === 'water' ? (active_chiller_fraction * (energy_it_yr / 1000) * 0.22) : 0.0;
const water_cost_yr = water_evap_m3 * waterPrice;

// Costo por fugas de refrigerantes
const total_fluid_charge = numRacks * 25; // 25kg por rack
const fluid_loss_yr = total_fluid_charge * (leakRate / 100);
const fluid_repl_cost_yr = fluid_loss_yr * fluidPrice;

const annual_opex = energy_sys_yr * tariff + fluid_repl_cost_yr + water_cost_yr;
const annual_opex_baseline = energy_baseline_yr * tariff;
const annual_opex_savings = annual_opex_baseline - annual_opex;

// Impacto de Carbono Neto (Scope 2 indirecto - Scope 1 directo de fugas)
const indirect_saved = (energy_saved_yr / 1000) * grid_factor;
const direct_leakage_emissions = (fluid_loss_yr / 1000) * fluidGwp;
const carbon_saved_yr = indirect_saved - direct_leakage_emissions;
\end{lstlisting}

\newpage

\section{Modelamiento Predictivo ARIMA (SQL - BigQuery ML)}
La siguiente consulta SQL es utilizada para entrenar el modelo de series de tiempo de temperatura exterior horaria directamente dentro de Google BigQuery:

\begin{lstlisting}[language=SQL]
-- 1. Entrenamiento del modelo predictivo multivariable
CREATE OR REPLACE MODEL `greenops_cooling.arima_temp_forecast`
OPTIONS(
  model_type='ARIMA_PLUS',
  time_series_timestamp_col='timestamp',
  time_series_data_col='temperature_c',
  time_series_id_col='location',
  auto_arima=TRUE,
  data_frequency='HOURLY',
  holiday_region='CL'
) AS
SELECT 
  CAST(timestamp AS TIMESTAMP) as timestamp, 
  temperature_c, 
  location
FROM `greenops_cooling.hourly_weather_data`
WHERE timestamp BETWEEN '2023-01-01 00:00:00' AND '2025-12-31 23:00:00';

-- 2. Ejecucion del pronostico a 24 horas con 95% de intervalo de confianza
SELECT
  location,
  forecast_time,
  forecast_value,
  standard_error,
  confidence_interval_lower_bound,
  confidence_interval_upper_bound
FROM
  ML.FORECAST(MODEL `greenops_cooling.arima_temp_forecast`, 
              STRUCT(24 AS horizon, 0.95 AS confidence_level));
\end{lstlisting}
